\documentclass[11pt]{article}
\usepackage[margin=1in]{geometry}
\usepackage{amsmath,amssymb,amsthm}
\usepackage{graphicx}
\usepackage[hidelinks]{hyperref}

\newtheorem{theorem}{Theorem}
\newtheorem{corollary}{Corollary}
\newtheorem{lemma}{Lemma}

\newcommand{\R}{\mathbb R}
\newcommand{\Sgn}{\operatorname{Sgn}}
\newcommand{\Tr}{\operatorname{Tr}}
\newcommand{\rank}{\operatorname{rank}}

\title{ETF Cases and a Rational Blow-Up Reduction for Basso's Orthogonal Projection Question}
\author{agent\_lane\_13}
\date{June 19, 2026}

\begin{document}
\maketitle

\section*{Status}

This packet records a likely valid partial result for Question 1.1 in
Giuliano Basso, \emph{Almost minimal orthogonal projections},
arXiv:2001.08698.  It does not settle the question for all dimensions.  It
proves the question for every dimension admitting a real maximal equiangular
tight frame, gives the concrete known cases $n=1,2,3,7,23$, and isolates a
rationality criterion that would turn Basso's almost-positive construction
into an exact construction.

\section*{Source Question}

Basso asks the following question in the introduction.

\begin{center}
\includegraphics[width=\textwidth]{figures/open_problem_crop.png}
\end{center}

Here $\Pi(E,F)$ denotes the relative projection constant, $\Pi_n$ denotes the
maximal real projection constant among $n$-dimensional Banach spaces, and the
orthogonal projection is taken with respect to the Euclidean structure on
$\R^d$.

\section*{Background from Basso}

For a symmetric sign matrix $S$ of size $N$ with diagonal entries equal to
$1$, let
\[
  \pi_n(S)=\lambda_1(S)+\cdots+\lambda_n(S),
\]
where the eigenvalues are ordered decreasingly.  Basso proves
\[
  \sup_{N\geq n}\sup_{S\in\mathcal S_N}\frac{1}{N}\pi_n(S)=\Pi_n.
\]
He also proves that if this supremum is attained by some finite sign matrix
$S\in\mathcal S_N$, then Question 1.1 has a positive answer in dimension $n$
\cite[Theorem 1.3 and Proposition 1.4]{Basso2001}.

\section*{Idea of the Proof}

There are two ways to make Basso's unweighted sign-matrix supremum attain its
value.

First, a real maximal equiangular tight frame produces a rank-$n$ orthogonal
projection $P$ whose entries have one diagonal magnitude and one off-diagonal
magnitude.  Hence $P$ is an affine function of its sign matrix $S$.  The top
$n$ eigenvalues of $S$ can then be computed exactly, and the 2022 theorem of
Deregowska and Lewandowska says that this value is the maximal projection
constant $\Pi_n$ whenever a real maximal ETF exists.

Second, Basso's proof already gives an exact weighted maximizer
$\sqrt D S\sqrt D$.  If the diagonal weights of $D$ are rational, replacing
the $i$-th vertex by a block of size proportional to $D_{ii}$ produces an
unweighted sign matrix whose nonzero spectrum is the spectrum of
$N\sqrt D S\sqrt D$.  Thus the unweighted supremum is attained exactly.

\section*{ETF Subcases}

\begin{theorem}[Maximal ETF criterion]
Suppose there is a real maximal equiangular tight frame in $\R^n$, that is,
an ETF with
\[
  N=\frac{n(n+1)}{2}
\]
vectors in $\R^n$.  Then Basso's Question 1.1 has a positive answer for this
integer $n$.
\end{theorem}

\begin{proof}
Let $x_1,\ldots,x_N$ be unit vectors forming a real ETF in $\R^n$, with
$N=n(n+1)/2$.  Put
\[
  U=\sqrt{n/N}\,[x_1\ \cdots\ x_N],
\]
so $UU^t=I_n$, and let $P=U^tU$.  Then $P$ is an $N\times N$ orthogonal
projection of rank $n$.  Set
\[
  a=\frac{n}{N},\qquad
  \alpha=\sqrt{ \frac{N-n}{n(N-1)} },\qquad
  b=a\alpha.
\]
The ETF property gives
\[
  p_{ii}=a,\qquad |p_{ij}|=b\quad (i\neq j).
\]
Let $S\in\mathcal S_N$ be the sign matrix of $P$, with $s_{ii}=1$ and
$s_{ij}=\operatorname{sign}(p_{ij})$ for $i\neq j$.  Then
\[
  P=bS+(a-b)I_N.
\]
Since $P$ has eigenvalue $1$ with multiplicity $n$ and eigenvalue $0$ with
multiplicity $N-n$, the matrix $S$ has eigenvalue
\[
  \frac{1-a+b}{b}
\]
with multiplicity $n$, and eigenvalue
\[
  -\frac{a-b}{b}
\]
with multiplicity $N-n$.  Therefore
\[
  \frac{1}{N}\pi_n(S)
  =\frac{n}{N}\,\frac{1-a+b}{b}
  =\frac{n}{N}+\frac{N-n}{N\alpha}
  =\frac{n}{N}\left(1+\sqrt{\frac{(N-1)(N-n)}{n}}\right).
\]
For $N=n(n+1)/2$ this is
\[
  \frac{2}{n+1}\left(1+\frac{n-1}{2}\sqrt{n+2}\right).
\]
Deregowska and Lewandowska prove that the real maximal projection constant
$\lambda_{\R}(n)$ equals this number whenever a real maximal ETF exists
\cite[Theorem 2.5]{DeregowskaLewandowska2206}.  This is Basso's $\Pi_n$ in
the real notation.  Hence $S$ attains Basso's supremum
\[
  \sup_{M\geq n}\sup_{T\in\mathcal S_M}\frac{1}{M}\pi_n(T)=\Pi_n.
\]
Basso's Proposition 1.4 now supplies an $n$-dimensional subspace
$E\subset\ell_1^N$ whose orthogonal projection is minimal and satisfies
\[
  \|P_E\|=\Pi(E,\ell_1^N)=\Pi_n.
\]
This is precisely a positive answer to Question 1.1 for this $n$.
\end{proof}

\begin{corollary}[Concrete dimensions]
Basso's Question 1.1 has a positive answer for
\[
  n=1,2,3,7,23.
\]
\end{corollary}

\begin{proof}
The case $n=1$ is immediate from the identity projection on $\ell_1^1$.  The
real maximal ETF cases listed by Deregowska and Lewandowska are
$n=2,3,7,23$ \cite[after Theorem 2.5]{DeregowskaLewandowska2206}.  Applying
the preceding theorem proves the claim.
\end{proof}

\section*{Rational-Weight Reduction}

\begin{theorem}[Exact rational blow-up criterion]
Fix $n\geq1$.  Suppose there are $m\geq n$, a sign matrix
$S\in\mathcal S_m$, and a diagonal matrix
\[
  D=\operatorname{diag}(d_1,\ldots,d_m),\qquad d_i>0,\quad \sum_i d_i=1,
\]
such that every $d_i$ is rational and
\[
  \pi_n(\sqrt D\,S\,\sqrt D)=\Pi_n.
\]
Then Basso's Question 1.1 has a positive answer for this $n$.
\end{theorem}

\begin{proof}
Choose positive integers $k_1,\ldots,k_m$ and let $N=\sum_i k_i$ so that
$d_i=k_i/N$.  Let $\widetilde S\in\mathcal S_N$ be the block blow-up of $S$
in which vertex $i$ is replaced by a block of size $k_i$, and the
$(i,j)$-block is constantly equal to $s_{ij}$.

Let $H\subset\R^N$ be the subspace of vectors that are constant on each block,
and write
\[
  u_i=\frac{{\mathbf 1}_{\text{block }i}}{\sqrt{k_i}}\qquad(1\leq i\leq m).
\]
The vectors $u_i$ form an orthonormal basis of $H$.  The subspace $H$ is
invariant under $\widetilde S$, and the matrix of $\widetilde S|_H$ in this
basis is
\[
  \bigl(s_{ij}\sqrt{k_ik_j}\bigr)_{i,j}
  =N\sqrt D\,S\,\sqrt D.
\]
The orthogonal complement $H^\perp$ consists of vectors whose coordinates sum
to zero on each block, and $\widetilde S$ annihilates $H^\perp$.  Thus the
spectrum of $\widetilde S$ is the spectrum of
$N\sqrt D\,S\,\sqrt D$ together with additional zeros.

Consequently
\[
  \frac{1}{N}\pi_n(\widetilde S)
  \geq \pi_n(\sqrt D\,S\,\sqrt D)
  =\Pi_n.
\]
Basso's formula for the unweighted supremum gives the reverse inequality for
every sign matrix, so equality holds:
\[
  \frac{1}{N}\pi_n(\widetilde S)=\Pi_n.
\]
Therefore $\widetilde S$ attains Basso's supremum.  Proposition 1.4 of
\cite{Basso2001} again gives the required subspace of $\ell_1^N$ and a
minimal orthogonal projection of norm $\Pi_n$.
\end{proof}

\section*{Verification Notes}

No numerical computation is used.  The ETF calculation only uses the standard
Gram-matrix form of a tight ETF and the exact equality statement for maximal
projection constants from Deregowska--Lewandowska.  The rational blow-up
argument is the exact version of the approximation-by-rational-weights step in
Basso's proof of the almost-positive theorem.

The key point in both arguments is not merely that large projection constants
are attained, but that Basso's unweighted sign-matrix supremum is attained.
That is exactly the hypothesis needed to invoke Basso's Proposition 1.4.

Bounded novelty checks were performed against the run indexes, the parsed
local source of arXiv:2001.08698, exact title/question web searches, and
searches for later projection-constant and maximal-ETF papers.  No full
all-$n$ resolution of Question 1.1 surfaced.

\section*{Limitations}

This packet is not a full solution.  It proves Question 1.1 for the maximal
ETF dimensions and gives an exact rationality criterion for the general
weighted maximizer.  Basso's Lemma 3.7 guarantees a positive diagonal
maximizer, but not rational diagonal entries.  Without rational weights the
block blow-up gives only the approximation already exploited in the source
paper.

\section*{Human Review Recommendation}

Review two points.  First, check the affine relation
$P=bS+(a-b)I_N$ for the normalized ETF Gram projection and the resulting
eigenvalue computation.  Second, check that the block blow-up has kernel
$H^\perp$ and compression $N\sqrt D S\sqrt D$ on block-constant vectors.  If
both checks pass, the reductions to Basso's Proposition 1.4 are direct.

\begin{thebibliography}{9}

\bibitem{Basso2001}
G. Basso,
\emph{Almost minimal orthogonal projections},
arXiv:2001.08698, 2020.

\bibitem{DeregowskaLewandowska2206}
B. Deregowska and B. Lewandowska,
\emph{A simple proof of the Grunbaum conjecture},
arXiv:2206.09454, 2022.

\end{thebibliography}

\end{document}
